\documentclass[10pt,a4paper]{amsart}
\usepackage[margin=19mm]{geometry}
\usepackage{amsmath,amssymb,mathtools}
\usepackage[scr=rsfs,cal=euler]{mathalfa}
\usepackage{xfrac}
\usepackage[unicode,hidelinks,bookmarks=false]{hyperref}
\newcommand{\ie}{i.e.\ }
\newcommand{\cf}{cf.\ }
\newcommand{\resp}{respectively\ }
\newcommand{\Subsection}{\subsection}
\providecommand{\nobreakdash}{\nobreak\hbox{-}}
\setlength{\emergencystretch}{2em}
\multlinegap0pt
\usepackage{mathtools}
\usepackage{amssymb}
\usepackage{algorithm}\usepackage{algpseudocode}
\usepackage{xcolor}
\newtheorem{proposition}{Proposition}
\usepackage{cleveref}
\crefname{proposition}{Proposition}{Propositions}
\DeclareMathOperator{\diag}{diag}
\title{Convergence Analysis of Randomized Subspace Normalized SGD under Heavy-Tailed Noise}
\author{Gaku Omiya, Pierre-Louis Poirion, Akiko Takeda}
\date{Source: arXiv:2601.20399v2 (CC BY 4.0)}
\begin{document}
\maketitle

\noindent\emph{Source and licence.} Convergence Analysis of Randomized Subspace Normalized SGD under Heavy-Tailed Noise. Version arXiv:2601.20399v2 (CC BY 4.0), distributed by arXiv under the Creative Commons Attribution 4.0 International licence, http://creativecommons.org/licenses/by/4.0/, as recorded in the arXiv licence metadata for this exact version and shown on the abstract page https://arxiv.org/abs/2601.20399v2. Full licence notice: \texttt{licenses/arxiv-2601.20399-CC-BY-4.0.txt}.\par\smallskip

\noindent\textit{Editorial scope.} Counted unit: the article's proposition app:prop:secondterm (source0.tex lines 2049--2058). The article's complete original proof of it is delivered with it (source0.tex lines 2059--2280, one \texttt{proof} environment of 222 lines). No mathematics is written, repaired or re-derived: the statement and the proof are the article's own lines, in the article's own notation, and no retained line is substituted or re-flowed.  Retained uncounted as the source-local context the proof needs: lines 604--612.  Nothing was omitted from any retained span, so the package carries no omission record.  No retained line contains a citation, so the wrapper carries no bibliography and no bibliography entry.  No retained line contains a figure environment, an image-inclusion command or drawing code, so no image is delivered and the package declares no asset dependency.  Editorial formatting only, in addition: an AMS \texttt{amsart} preamble that restates the article's own declarations and macros used by the retained lines, namely the environments \texttt{proposition} on separate counters, together with the standard packages those lines need.  The retained environments print in this document as Proposition 1 in order of appearance, whereas the article prints the same items with the numbering of its own full document.  The complete native source of the article as distributed in the arXiv e-print for this exact version is bundled unchanged as \texttt{sources/193519/source0.tex}.
\begin{proposition}\label{prop:secondterm}
Let $S$ be a $d\times d$ square matrix, and $P$ be a Haar matrix. Then,
\[
\mathbb{E}\left[\frac{v^\top PP^\top S PP^\top v}{\|P^\top v\|^2}\right]
=
\frac{d(r-1)}{r(d-1)}\frac{v^\top Sv}{\|v\|^2}
+\frac{d-r}{r(d-1)}\operatorname{tr}(S).
\]
\end{proposition}

\begin{proposition}[Restatement of \Cref{prop:secondterm}]
\label{app:prop:secondterm}
Let $S$ be a $d\times d$ square matrix, and $P$ be a Haar matrix. Then,
\[
\mathbb{E}\left[\frac{v^\top PP^\top S PP^\top v}{\|P^\top v\|^2}\right]
=
\frac{d(r-1)}{r(d-1)}\frac{v^\top Sv}{\|v\|^2}
+\frac{d-r}{r(d-1)}\operatorname{tr}(S).
\]
\end{proposition}

\begin{proof}
Fix a vector $v \in \mathbb{R}^d$.  
There exists an orthogonal matrix $\hat Q \in O(d)$ such that
\[
v = \|v\| \hat Q e_1,
\]
where $e_1 \in \mathbb{R}^d$ denotes the vector whose first component is $1$ and whose remaining components are $0$.

By the invariance of the Haar matrix, $\hat Q P$ has the same distribution as $P$.  
Define
\[
\hat S \coloneqq \hat Q^\top S \hat Q .
\]
Then we have
\begin{align}
\mathbb{E}\!\left[\frac{v^\top PP^\top S PP^\top v}{\|P^\top v\|^2}\right]
&= \mathbb{E}\!\left[\frac{v^\top \hat Q PP^\top \hat Q^\top S \hat Q PP^\top \hat Q^\top v}{\|P^\top \hat Q^\top v\|^2}\right] \notag\\
&= \mathbb{E}\!\left[\frac{\|v\| e_1^\top PP^\top \hat Q^\top S \hat Q PP^\top \|v\| e_1}{\|P^\top \|v\| e_1\|^2}\right] \notag\\
&= \mathbb{E}\!\left[\frac{e_1^\top PP^\top \hat S PP^\top e_1}{\|P^\top e_1\|^2}\right].
\label{eq:LHS}
\end{align}

Define the random vector
\[
u \coloneqq \frac{PP^\top e_1}{\|P^\top e_1\|}.
\]
Then
\begin{align}
\mathbb{E}\!\left[\frac{e_1^\top PP^\top \hat S PP^\top e_1}{\|P^\top e_1\|^2}\right]
&= \mathbb{E}[u^\top \hat S u] \notag\\
&= \operatorname{tr}\!\left(\mathbb{E}[uu^\top]\hat S\right),
\label{eq:RHS}
\end{align}
where we used the identity
$\mathbb{E}[u^\top \hat S u]
= \mathbb{E}[\operatorname{tr}(u^\top \hat S u)]
= \mathbb{E}[\operatorname{tr}(uu^\top \hat S)]
= \operatorname{tr}(\mathbb{E}[uu^\top]\hat S)$.

Combining \eqref{eq:LHS} and \eqref{eq:RHS}, we obtain
\[
\mathbb{E}\!\left[\frac{v^\top PP^\top S PP^\top v}{\|P^\top v\|^2}\right]
= \operatorname{tr}(M\hat S),
\quad
M \coloneqq \mathbb{E}[uu^\top].
\]

We now characterize the structure of $M$.
Consider the subset
\[
\mathcal{S} \coloneqq \{Q \in O(d) \mid Q e_1 = e_1\}.
\]
For any $Q \in \mathcal{S}$, by the invariance of the Haar matrix,
\begin{align*}
Q^\top M Q
&= \mathbb{E}[Q^\top uu^\top Q] \\
&= \mathbb{E}\!\left[\frac{Q^\top PP^\top e_1 e_1^\top PP^\top Q}{\|P^\top e_1\|^2}\right] \\
&= \mathbb{E}\!\left[\frac{Q^\top PP^\top Q e_1 e_1^\top Q^\top PP^\top Q}{\|P^\top Q e_1\|^2}\right] \\
&= \mathbb{E}\!\left[\frac{(Q^\top P)(Q^\top P)^\top e_1 e_1^\top (Q^\top P)(Q^\top P)^\top}{\|(Q^\top P)^\top e_1\|^2}\right] \\
&= \mathbb{E}\!\left[\frac{P P^\top e_1 e_1^\top PP^\top}{\|P^\top e_1\|^2}\right] \\
&= \mathbb{E}[uu^\top] \\
&= M.
\end{align*}

This invariance implies strong structural constraints on $M$.
First, fix $i \in \{2,\dots,d\}$ and define the matrix
\[
A_i \coloneqq \diag(1,\dots,1,-1,1,\dots,1) \in O(d-1),
\]
where the $(i-1)$-th diagonal entry is $-1$ and the remaining $d-2$ diagonal
entries are equal to $1$.
Define
\[
Q \coloneqq
\begin{pmatrix}
1 & \mathbf{0} \\
\mathbf{0} & A_i
\end{pmatrix}.
\]
Then $Q \in \mathcal{S}$ by construction.

Since $Q$ has $-1$ only at the $i$-th diagonal entry and $1$ at all other
diagonal entries, we compute $Q^\top M Q$ explicitly.
Writing $M=(m_{jk})_{1\leq j,k\leq d}$, we have
\begin{align*}
M
= Q^\top M Q
= \begin{pmatrix}
  m_{1,1} & \cdots & m_{1,i-1} & -m_{1,i} & m_{1,i+1} & \cdots & m_{1,d} \\[4pt]
  \vdots  &        & \vdots    & \vdots   & \vdots    &        & \vdots \\[4pt]
  m_{i-1,1} & \cdots & m_{i-1,i-1} & -m_{i-1,i} & m_{i-1,i+1} & \cdots & m_{i-1,d} \\[6pt]
  -m_{i,1} & \cdots & -m_{i,i-1} & m_{i,i} & -m_{i,i+1} & \cdots & -m_{i,d} \\[6pt]
  m_{i+1,1} & \cdots & m_{i+1,i-1} & -m_{i+1,i} & m_{i+1,i+1} & \cdots & m_{i+1,d} \\[4pt]
  \vdots &        & \vdots & \vdots & \vdots &        & \vdots \\[4pt]
  m_{d,1} & \cdots & m_{d,i-1} & -m_{d,i} & m_{d,i+1} & \cdots & m_{d,d}
\end{pmatrix}.
\end{align*}
Since this equality holds for all $i \in \{2,\dots,d\}$, it follows that
\[
m_{jk}=0 \quad \text{for all } j\neq k,
\]
and hence $M$ must be a diagonal matrix.
Therefore, $M$ can be written as
\[
M=
\begin{pmatrix}
m_{1,1} & 0 & \cdots & 0 \\
0 & m_{2,2} & \cdots & 0 \\
\vdots & \vdots & \ddots & \vdots \\
0 & 0 & \cdots & m_{d,d}
\end{pmatrix}.
\]

Next, take $i,j \in \{2,\dots,d\}$ with $i \le j$, and let $Q$ be the permutation
matrix corresponding to the transposition $(i,j)$, i.e.,
\[
Q\coloneqq
\begin{pmatrix}
1 &        &        &        &        &        &        \\
  & \ddots &        &        &        &        &        \\
  &        & 0      &        & 1      &        &        \\
  &        &        & \ddots &        &        &        \\
  &        & 1      &        & 0      &        &        \\
  &        &        &        &        & \ddots &        \\
  &        &        &        &        &        & 1
\end{pmatrix}.
\]
The lower-right block of $Q$ is a $(d-1)\times(d-1)$ permutation matrix, and
hence belongs to $O(d-1)$.
Therefore, $Q \in \mathcal{S}$.

Using again the invariance $Q^\top M Q=M$, we obtain
\begin{align*}
M
&= Q^\top M Q \\
&= \diag(
m_{1,1},\dots,m_{i-1,i-1},
m_{j,j},
m_{i+1,i+1},\dots,m_{j-1,j-1},
m_{i,i},
m_{j+1,j+1},\dots,m_{d,d}
).
\end{align*}
Since this holds for arbitrary $i,j \in \{2,\dots,d\}$ with $i \le j$, we conclude
that
\[
m_{2,2}=m_{3,3}=\cdots=m_{d,d}.
\]

Defining $\alpha \coloneqq m_{1,1}$ and $\beta \coloneqq m_{2,2}$, we finally
obtain
\[
M=
\begin{pmatrix}
\alpha & 0 & \cdots & 0 \\
0 & \beta & \cdots & 0 \\
\vdots & \vdots & \ddots & \vdots \\
0 & 0 & \cdots & \beta
\end{pmatrix}
= (\alpha-\beta)e_1e_1^\top + \beta I_d.
\]

We now determine $\alpha$ and $\beta$.
First,
\begin{align*}
e_1^\top M e_1
&= \mathbb{E}[e_1^\top uu^\top e_1] \\
&= \mathbb{E}\!\left[e_1^\top \frac{PP^\top e_1 e_1^\top PP^\top}{\|P^\top e_1\|^2} e_1\right] \\
&= \mathbb{E}[\|P^\top e_1\|^2] \\
&= e_1^\top \mathbb{E}[PP^\top] e_1 \\
&= 1.
\end{align*}
On the other hand,
\[
e_1^\top M e_1 = \alpha,
\]
and hence $\alpha=1$.

Next, we compute the trace of $M$:
\begin{align*}
\operatorname{tr}(M)
&= \operatorname{tr}(\mathbb{E}[uu^\top]) \\
&= \mathbb{E}[\operatorname{tr}(uu^\top)] \\
&= \mathbb{E}[\|u\|^2] \\
&= \mathbb{E}\!\left[\frac{e_1^\top PP^\top PP^\top e_1}{\|P^\top e_1\|^2}\right] \\
&= \frac{d}{r}.
\end{align*}
On the other hand,
\begin{align*}
\operatorname{tr}(M)
&= \operatorname{tr}((1-\beta)e_1e_1^\top + \beta I_d) \\
&= (1-\beta) + \beta d.
\end{align*}
Equating the two expressions yields
\[
(d-1)\beta = \frac{d}{r} - 1,
\quad
\beta = \frac{d-r}{r(d-1)}.
\]

Therefore,
\[
M
= \frac{d(r-1)}{r(d-1)} e_1e_1^\top
+ \frac{d-r}{r(d-1)} I_d.
\]

Finally,
\begin{align*}
\mathbb{E}\!\left[\frac{v^\top PP^\top S PP^\top v}{\|P^\top v\|^2}\right]
&= \operatorname{tr}(M\hat S) \\
&= \operatorname{tr}\!\left(
\frac{d(r-1)}{r(d-1)} e_1e_1^\top \hat S
+ \frac{d-r}{r(d-1)} \hat S
\right) \\
&= \frac{d(r-1)}{r(d-1)} e_1^\top \hat S e_1
+ \frac{d-r}{r(d-1)} \operatorname{tr}(\hat S) \\
&= \frac{d(r-1)}{r(d-1)} \frac{v^\top S v}{\|v\|^2}
+ \frac{d-r}{r(d-1)} \operatorname{tr}(S),
\end{align*}
which completes the proof.
\end{proof}

\end{document}
